\chapter{Hom 细节全拆：指纹怎么滚、三边怎么改、如何调弦}
\label{ch:hom-details}

\section{Gear 表：房间专属骰子}

\begin{metaphor}[赌场换骰]
\texttt{table\_seed} 决定 256 面``骰子表'' $G[0..255]$。\\
Hom 房间的骰子与 FastCDC / G-TCDC \textbf{不是同一盒} —— 隔离，免得铃声模式撞车。
\end{metaphor}

真名：\texttt{InitKeyedGearTable(topo\_table\_seed)}。

\section{指纹滚动：传送带公式}

窗宽 $w$。每进一步：左边字节贡献掉出，右边加进来：
\begin{equation}
h \leftarrow 2h + G[b_{\mathrm{in}}] - G[b_{\mathrm{out}}].
\end{equation}

\begin{figure}[htbp]
\centering
\begin{tikzpicture}[font=\small]
  \foreach \i in {0,...,5} {
    \node[bytecell] (b\i) at (\i*0.9,0) {$b$};
  }
  \draw[thick, Gen3Color] (-0.3,-0.5) rectangle (5.1,0.5);
  \node[font=\footnotesize, CutRed, above] at (0,0.7) {out};
  \node[font=\footnotesize, OkGreen, above] at (5.4,0.7) {in};
  \draw[-{Stealth}, CutRed] (-0.8,0)--(-0.35,0);
  \draw[-{Stealth}, OkGreen] (5.15,0)--(6.0,0);
  \node[softbox=Gen3Color, below=0.9cm of b2] {$h$ 更新};
\end{tikzpicture}
\caption{滚筒只关心进出两颗珠对指纹的加减。}
\label{fig:gear-roll-detail}
\end{figure}

\textbf{叮？} $(h \mathbin{\&} M)=0$。$M$ 是一串低位 1 的掩码（越松越好叮）。

\section{颜色键}

珠串颜色不是原始字节，而是 $k=G[b]\mathbin{\&}0xFF$（取骰子值低 8 位）。\\
同字节在不同 seed 下可能不同色。

\section{三边捷径（PrimaryTopoDelta）逐项}

滑动：丢掉最左，右侧进入 $k_{\mathrm{in}}$。只有这几根边的``在/不在''可能变：

\begin{figure}[htbp]
\centering
\begin{tikzpicture}[font=\scriptsize]
  \foreach \i/\lab in {0/k0, 1/k1, 2/mid, 3/kw-2, 4/kw-1} {
    \node[circle,draw,fill=Gen3Color!20,minimum size=0.7cm] (p\i) at (\i*1.35,0) {\lab};
  }
  \node[circle,draw,fill=orange!40,minimum size=0.7cm] (kin) at (7.2,0) {kin};
  \draw[dashed, CutRed, thick] (p0)--(p1);
  \node[font=\tiny, CutRed, above] at (0.65,0.55) {可能消失};
  \draw[dashed, CutRed, thick] (p3)--(p4);
  \node[font=\tiny, CutRed, above] at (4.7,0.55) {可能消失};
  \draw[dashed, OkGreen, thick] (p3)--(kin);
  \draw[dashed, OkGreen, thick] (kin) to[bend left=25] (p0);
  \node[font=\tiny, OkGreen, below] at (6.2,-0.7) {可能新生};
\end{tikzpicture}
\caption{O(1) 只改最多三/四次``异色边计数''，不必重扫全窗。}
\label{fig:three-edges}
\end{figure}

\begin{toybox}[数字]
旧边数 $\mathrm{ed}=4$。运算后 $\mathrm{ed}=5$。$\Delta=+1\neq0$ $\Rightarrow$ topoOK。\\
连通分量口语：$C=1+\mathrm{ed}$（整条路径图）。
\end{toybox}

\section{lazy：没叮就不数珠}

\begin{figure}[htbp]
\centering
\begin{tikzpicture}[font=\small]
  \node[softbox=gray] (a) {1000 步};
  \node[softbox=Gen3Color, right=0.8cm of a] (b) {约 avg 分之一叮};
  \node[softbox=Gen3Color, right=0.8cm of b] (c) {才 $O(w)$ 数珠};
  \draw[-{Stealth}, thick] (a)--(b)--(c);
\end{tikzpicture}
\caption{这就是 Hom 能接近 Stream 的原因。}
\label{fig:lazy-amortize}
\end{figure}

对照测验：\texttt{SlideViaRebuild} 全量 UF $=$ 黄金答案；生产用三边捷径，单测要求一致。

\section{调弦两步曲}

\begin{enumerate}
  \item 临时用粗 mask 扫标定样本，统计``叮的时候珠串也变''的比例 $P$；
  \item 令有效平均 $\approx \mathrm{avg}\cdot P$，再 \texttt{BuildMask} 得生产 $M$；
  \item 把 $1000P$ 写成 \texttt{calib\_permille} 存进超块。
\end{enumerate}

\begin{figure}[htbp]
\centering
\begin{tikzpicture}[font=\small]
  \node[softbox=Gen3Color] (1) {弹 1MB 样音};
  \node[softbox=Gen3Color, right=0.5cm of 1] (2) {读 $P$};
  \node[softbox=Gen3Color, right=0.5cm of 2] (3) {拧 mask};
  \node[softbox=OkGreen, right=0.5cm of 3] (4) {平均块长合格};
  \draw[-{Stealth}, thick] (1)--(2)--(3)--(4);
\end{tikzpicture}
\caption{旧仓 \texttt{calib=0}：退回 \texttt{avg >> topo\_shift} 老旋钮。}
\label{fig:hom-calib-pipe}
\end{figure}

\section{流式小条：UF 不必装箱带走}

Resume 主要带走：指纹 $h$、扫到哪里。珠串 UF \textbf{可不装箱}——下次叮了再现场重串珠。\\
整文件切点 $=$ 多个 512KB 盒子拼起来的切点（parity）。

\section{AVX2：一次问 8 格}

像八车道同时探头，但刀口规则不变 —— 加速外壳，不改判决。

\section{细节对照表}

\begin{center}
\small
\begin{tabular}{@{}lp{7.5cm}@{}}
\toprule
技术点 & 形象说法 \\
\midrule
\texttt{InitWindowHash} & 开窗时先灌满滚筒 \\
\texttt{BuildTopoMasks} & 按 $P$ 或 shift 选铃灵敏度 \\
\texttt{ScanHomCut} & 在 min..max 里找第一刀 \\
\texttt{ProcessHomChunk} & 一块香肠内的扫描状态机 \\
\texttt{TopoGearParityTest} & 整条 vs 分段必须同刀 \\
\bottomrule
\end{tabular}
\end{center}
